\documentclass[11pt,a4paper]{article}
\usepackage[T1]{fontenc}
\usepackage[utf8]{inputenc}
\usepackage{amsmath,amssymb,amsthm}
\usepackage[margin=1in]{geometry}
\usepackage[colorlinks=true,linkcolor=blue,urlcolor=blue]{hyperref}
\theoremstyle{plain}
\newtheorem{theorem}{Theorem}
\newtheorem{lemma}[theorem]{Lemma}
\newtheorem{proposition}[theorem]{Proposition}
\newtheorem{corollary}[theorem]{Corollary}
\theoremstyle{definition}
\newtheorem{remark}[theorem]{Remark}
% A quoted conjecture can be a single hundred-term formula whose terms are thirty-digit
% numbers. TeX cannot hyphenate those, so without a little stretch every such quote runs off
% the page; the linked-a-file papers are all of that shape.
\setlength{\emergencystretch}{3em}

\title{An exact count for OEIS A228462: windowed maxima over a growing alphabet}
\author{Adrian Perez Fontelles\\ \small Independent researcher}
\date{15 September 2026}
\begin{document}
\maketitle

\begin{abstract}
OEIS A228462 counts the arrays obtainable as the maxima of $3$ adjacent elements of a length
$7$ array with entries in $\{0,\dots,n\}$, and carries an empirical polynomial for that
count. The count is not a transfer-matrix count: the length is fixed and it is the ALPHABET
that grows with $n$. It is nevertheless exactly computable, because membership in the image
depends only on the relative order of the entries of the candidate array and not on their
values, so the count is a fixed non-negative combination of binomial coefficients
$\binom{n+1}{m}$. Computing that combination settles the empirical claim as an identity in
$\mathbb{Q}[n]$.
\end{abstract}

\noindent\small 2020 Mathematics Subject Classification. 05A15, 05A05, 06A07.\normalsize

\section{The conjecture}

The entry's first terms are $17, 91, 310, 821, 1847, 3703, 6812, 11721,\dots$ and it states

\begin{quote}\raggedright Empirical: a(n) = (2/15)*n\^{}5 + (7/6)*n\^{}4 + (25/6)*n\^{}3 + (19/3)*n\^{}2 + (21/5)*n + 1.\end{quote}

\section{The object}

Fix $k=7$ and $w=3$, and put $L=k-w+1=5$. For $a=(a_1,\dots,a_k)\in\{0,\dots,n\}^k$
let
\[
  \Phi(a)=\bigl(\max(a_1,\dots,a_w),\ \max(a_2,\dots,a_{w+1}),\ \dots,\
  \max(a_{k-w+1},\dots,a_k)\bigr)\in\{0,\dots,n\}^{L} .
\]
The entry counts $\lvert\Phi(\{0,\dots,n\}^k)\rvert$.

\section{Membership in the image is decided by the order type}

Write $W_j=\{j,j+1,\dots,j+w-1\}$ for the $j$-th window, $1\le j\le L$.

\begin{lemma}\label{lem:greedy}
Let $b\in\{0,\dots,n\}^{L}$ and define $\hat a\in\{0,\dots,n\}^k$ by
$\hat a_i=\min\{\,b_j : i\in W_j\,\}$. Then $b\in\Phi(\{0,\dots,n\}^k)$ if and only if
$\Phi(\hat a)=b$.
\end{lemma}

\begin{proof}
Each $\hat a_i$ is one of the $b_j$, so $\hat a$ does lie in $\{0,\dots,n\}^k$ and the
condition is meaningful. If $\Phi(\hat a)=b$ there is nothing to prove. Conversely suppose
$\Phi(a)=b$ for some $a$. For every $i\in W_j$ we have $a_i\le\max_{l\in W_j}a_l=b_j$, so
$a_i\le\min\{b_j : i\in W_j\}=\hat a_i$: the array $\hat a$ dominates every witness
componentwise. Hence
$\max_{i\in W_j}\hat a_i\ge\max_{i\in W_j}a_i=b_j$. The reverse inequality is immediate
from the definition, since $\hat a_i\le b_j$ for each $i\in W_j$. So $\Phi(\hat a)=b$.
\end{proof}

Both the construction of $\hat a$ and the test $\Phi(\hat a)=b$ use only comparisons between
entries of $b$. Consequently:

\begin{corollary}\label{cor:ordertype}
Whether $b$ lies in the image depends only on the \emph{order type} of $b$ -- the weak
ordering of its $L$ coordinates -- and not on $n$ or on the values themselves.
\end{corollary}

\begin{proof}
Let $\sigma$ be any strictly increasing map defined on the values of $b$. Applying $\sigma$
entrywise commutes with $\min$ and $\max$, so it carries $\hat a$ for $b$ to $\hat a$ for
$\sigma(b)$ and preserves the equalities tested in Lemma~\ref{lem:greedy}. Any two arrays with
the same order type differ by such a $\sigma$.
\end{proof}

\section{The count}

For $1\le m\le L$ let $A_m$ be the number of order types of $L$ coordinates that use exactly
$m$ distinct values and are achievable. An array $b\in\{0,\dots,n\}^{L}$ using exactly $m$
distinct values is determined by the $m$-element subset of $\{0,\dots,n\}$ it uses together
with its order type, and every pair occurs exactly once; so by Corollary~\ref{cor:ordertype},
\[
  \lvert\Phi(\{0,\dots,n\}^k)\rvert=\sum_{m=1}^{5}A_m\binom{n+1}{m}
  \qquad\text{for every } n\ge 0 .
\]
Enumerating the weak orderings of $5$ elements and applying Lemma~\ref{lem:greedy} to each
gives $A_{1}=1, A_{2}=15, A_{3}=43, A_{4}=44, A_{5}=16$, so
\[
  \lvert\Phi(\{0,\dots,n\}^k)\rvert=1\binom{n+1}{1} + 15\binom{n+1}{2} + 43\binom{n+1}{3} + 44\binom{n+1}{4} + 16\binom{n+1}{5}=\frac{2 n^{5}}{15} + \frac{7 n^{4}}{6} + \frac{25 n^{3}}{6} + \frac{19 n^{2}}{3} + \frac{21 n}{5} + 1 .
\]

\section{Conclusion}

The right-hand side is the polynomial the entry states, so the empirical claim of OEIS A228462 is
true for every $n\ge 0$. The argument is finite and exact: it enumerates the weak orderings of
$5$ elements, of which there are 119 achievable, and performs no fitting to the
published terms. Those terms are reproduced by the formula as a check, not as evidence.

\end{document}
